\documentclass[11pt, a4paper]{article}
\usepackage[utf8]{inputenc}
\usepackage[french]{babel}
\usepackage{lmodern}
\usepackage[T1]{fontenc}
\usepackage{geometry}
\geometry{margin=2.5cm}
\usepackage{hyperref}
\hypersetup{
    colorlinks=true,
    linkcolor=blue,
    urlcolor=blue,
    bookmarksnumbered=true,
    pdftitle={Manuel d'Utilisation Pangol1},
    pdfauthor={Pangol1 Team}
}
\usepackage{enumitem}
\usepackage{graphicx}
\usepackage{xcolor}

\title{Manuel d'Utilisation : Pangol1}
\author{}
\date{}

\addto\captionsfrench{\renewcommand{\contentsname}{Sommaire}}

\begin{document}

\maketitle

\section*{À propos de ce document}
\addcontentsline{toc}{section}{À propos de ce document}
Ce manuel fournit une documentation complète sur l'utilisation de \textbf{Pangol1}. Il couvre tous les aspects de l'application, de l'importation de fichiers CSV à la génération de graphiques SVG.

\vspace{2ex}
\noindent
{\small \textit{Dernière mise à jour : \today}}

\newpage

% Sommaire manuscrit
\section*{Sommaire}
\addcontentsline{toc}{section}{Sommaire}

\begin{enumerate}
    \item Introduction \\
    \quad 1.1 Qu'est-ce que Pangol1 ? \\
    \quad 1.2 Fonctionnalités Principales
    
    \item Démarrage Rapide \\
    \quad 2.1 Lancer l'Application \\
    \quad 2.2 Première Utilisation
    
    \item Interface Utilisateur \\
    \quad 3.1 Organisation Générale \\
    \quad 3.2 Barre de Menus
    
    \item Chargement des Données \\
    \quad 4.1 Importer un Fichier CSV \\
    \quad 4.2 Format des Fichiers \\
    \quad 4.3 Changement de Langue
    
    \item Exploration des Données \\
    \quad 5.1 Table Interactive \\
    \quad 5.2 Tri des Données \\
    \quad 5.3 Filtrage des Données \\
    \quad 5.4 Masquage de Colonnes
    
    \item Génération de Graphiques \\
    \quad 6.1 Types de Graphiques Disponibles \\
    \quad 6.2 Configuration du Graphique \\
    \quad 6.3 Personnalisation Visuelle
    
    \item Fonctionnalités Supplémentaires \\
    \quad 7.1 Mode Nuit \\
    \quad 7.2 Fichiers Récents
    
    \item Résolution de Problèmes \\
    \quad 8.1 Problèmes Courants \\
    \quad 8.2 Conseils d'Utilisation
    
    \item Support \\
    \quad 9.1 Obtenir de l'Aide \\
    \quad 9.2 Documentation Technique
\end{enumerate}

\newpage

\section{Introduction}

\subsection{Qu'est-ce que Pangol1 ?}
\textbf{Pangol1} est une application graphique permettant de générer des graphiques SVG à partir de fichiers CSV. Elle offre une interface intuitive pour charger, explorer et visualiser vos données sous forme de graphiques professionnels.

\subsection{Fonctionnalités Principales}
\begin{itemize}
    \item Importation de fichiers CSV depuis votre ordinateur
    \item Exploration interactive des données via une table
    \item Génération de trois types de graphiques : Camembert, Barres, Colonnes
    \item Personnalisation complète de l'apparence des graphiques
    \item Export des graphiques en format SVG
    \item Support de 9 langues
    \item Mode nuit automatique (18h à 6h)
\end{itemize}

\section{Démarrage Rapide}

\subsection{Lancer l'Application}
Pangol1 fonctionne sur Windows, Linux et macOS. Pour lancer l'application :

\begin{enumerate}
    \item Lancez le fichier \textbf{Pangol1.jar} ou exécutez la commande : \texttt{java -jar Pangol1.jar}
    \item L'interface graphique s'affiche
    \item Vous pouvez maintenant charger vos données et créer des graphiques
\end{enumerate}

\subsection{Première Utilisation}
\begin{enumerate}
    \item Lancez Pangol1
    \item Allez dans \textbf{File} → \textbf{Open} pour charger un fichier CSV
    \item Sélectionnez un fichier CSV sur votre ordinateur
    \item Les données s'affichent automatiquement dans la table à droite
    \item Configurez votre graphique via le panneau de configuration à gauche
    \item Exportez votre graphique en SVG si désiré
\end{enumerate}

\section{Interface Utilisateur}

\subsection{Organisation Générale}
L'interface de Pangol1 est organisée en trois zones principales :

\begin{itemize}
    \item \textbf{Panneau de Configuration (Gauche)} : Permet de configurer le graphique (type, données, apparence)
    
    \item \textbf{Zone de Rendu (Centre)} : Affiche le graphique généré en temps réel
    
    \item \textbf{Table de Données (Droite)} : Affiche les lignes du fichier CSV importé avec système de pagination
\end{itemize}

\subsection{Barre de Menus}
La barre de menus contient les options suivantes :

\begin{itemize}
    \item \textbf{File} : Gestion des fichiers (Open, Save, Export, Exit)
    \item \textbf{Help} : Aide et paramètres (Languages, About)
\end{itemize}

\section{Chargement des Données}

\subsection{Importer un Fichier CSV}
Pour charger un fichier CSV :

\begin{enumerate}
    \item Cliquez sur \textbf{File} → \textbf{Open}
    \item Une fenêtre de sélection s'ouvre
    \item Naviguez et sélectionnez votre fichier CSV
    \item Cliquez sur \textbf{Ouvrir}
    \item Le fichier est chargé et les données s'affichent dans la table
\end{enumerate}

\subsection{Format des Fichiers}
Pangol1 accepte les fichiers au format \textbf{CSV} (Comma-Separated Values) :

\begin{itemize}
    \item Les colonnes sont séparées par des virgules ou des points-virgules
    \item La première ligne contient généralement les en-têtes de colonnes
    \item Les fichiers doivent être encodés en UTF-8 pour un fonctionnement optimal
\end{itemize}

\subsection{Changement de Langue}
Pour changer la langue de l'interface :

\begin{enumerate}
    \item Cliquez sur \textbf{Help} → \textbf{Languages}
    \item Sélectionnez votre langue parmi les 9 disponibles :
    \begin{itemize}
        \item Français (défaut)
        \item Anglais
        \item Chinois simplifié
        \item Espagnol
        \item Arabe
        \item Portugais
        \item Russe
        \item Japonais
        \item Indonésien
    \end{itemize}
    \item L'interface se met à jour immédiatement
\end{enumerate}

\section{Exploration des Données}

\subsection{Table Interactive}
La table de données à droite affiche le contenu de votre fichier CSV :

\begin{itemize}
    \item \textbf{Pagination} : Utilisez les boutons \textbf{First, Previous, Next, Last} pour naviguer entre les pages
    \item \textbf{En-têtes de Colonnes} : Cliquez sur les en-têtes pour accéder aux options de tri et filtrage
    \item Chaque ligne représente un enregistrement du fichier
\end{itemize}

\subsection{Tri des Données}
Pour trier une colonne :

\begin{enumerate}
    \item Cliquez sur l'en-tête de la colonne
    \item Un menu contextuel apparaît
    \item Choisissez \textbf{Ascendant} ou \textbf{Descendant}
    \item Les données se réorganisent automatiquement
\end{enumerate}

\subsection{Filtrage des Données}
Pour filtrer les données :

\begin{enumerate}
    \item Cliquez sur l'en-tête de la colonne
    \item Sélectionnez \textbf{Top N} pour afficher les N valeurs les plus hautes
    \item Ou sélectionnez \textbf{Bottom N} pour afficher les N valeurs les plus basses
    \item Seules les lignes correspondantes s'affichent dans la table
\end{enumerate}

\subsection{Masquage de Colonnes}
Pour masquer une colonne :

\begin{enumerate}
    \item Cliquez sur l'en-tête de la colonne
    \item Sélectionnez \textbf{Hide} ou \textbf{Masquer}
    \item La colonne disparaît de la table
\end{enumerate}

\section{Génération de Graphiques}

\subsection{Types de Graphiques Disponibles}
Pangol1 propose trois types de graphiques :

\subsubsection{Camembert (Pie Chart)}
\begin{itemize}
    \item Utilisation : Visualiser les proportions et les parts
    \item Cas d'usage : Répartition de pourcentages, composition d'un ensemble
    \item Recommandation : Optimal avec 3 à 6 catégories
\end{itemize}

\subsubsection{Barres Horizontales (Bar Chart)}
\begin{itemize}
    \item Utilisation : Comparer des valeurs entre catégories
    \item Cas d'usage : Comparaisons, classements, benchmarking
    \item Recommandation : Adapté aux longs labels
\end{itemize}

\subsubsection{Colonnes Verticales (Columns Chart)}
\begin{itemize}
    \item Utilisation : Afficher l'évolution temporelle ou les variations
    \item Cas d'usage : Série temporelle, tendances, comparaisons verticales
    \item Recommandation : Standard pour l'analyse temporelle
\end{itemize}

\subsection{Configuration du Graphique}

\subsubsection{Sélection du Type}
Dans le panneau de configuration :

\begin{enumerate}
    \item Localisez la section \textbf{Type}
    \item Sélectionnez l'un des trois types de graphiques
    \item Le graphique se met à jour automatiquement
\end{enumerate}

\subsubsection{Sélection des Données}
\begin{enumerate}
    \item Sélectionnez la colonne pour l'axe X (catégories)
    \item Sélectionnez la colonne pour l'axe Y (valeurs numériques)
    \item Le graphique se régénère avec les nouvelles données
\end{enumerate}

\subsection{Personnalisation Visuelle}

\subsubsection{Couleurs}
\begin{itemize}
    \item Dans la section \textbf{Design}, sélectionnez une couleur de remplissage
    \item Chaque graphique a sa propre palette de couleurs
    \item Les modifications s'appliquent immédiatement
\end{itemize}

\subsubsection{Transparence (Opacity)}
\begin{itemize}
    \item Ajustez le niveau de transparence des éléments graphiques
    \item Valeurs de 0\% (invisible) à 100\% (opaque)
    \item Utile pour visualiser les données qui se chevauchent
\end{itemize}

\subsubsection{Épaisseur du Trait (Stroke Width)}
\begin{itemize}
    \item Contrôle l'épaisseur des bordures des éléments
    \item Valeurs recommandées : 1 à 5 pixels
    \item Augmentez pour une meilleure visibilité sur petit écran
\end{itemize}

\subsubsection{Labels des Axes}
\begin{itemize}
    \item \textbf{X Label} : Titre de l'axe horizontal
    \item \textbf{Y Label} : Titre de l'axe vertical
    \item Ces labels s'affichent sur le graphique généré
\end{itemize}

\subsubsection{Nombre de Graduations}
\begin{itemize}
    \item \textbf{Tick Count} : Nombre de marques sur les axes
    \item Augmentez pour plus de détails
    \item Diminuez pour une meilleure lisibilité
\end{itemize}

\subsubsection{Échelle du Graphique}
\begin{itemize}
    \item \textbf{Linéaire} : Échelle proportionnelle standard
    \item \textbf{Logarithmique} : Pour les données avec grande disparité
    \item La scale logarithmique améliore la lisibilité quand certaines valeurs dominent d'autres
\end{itemize}

\section{Fonctionnalités Supplémentaires}

\subsection{Mode Nuit}
\begin{itemize}
    \item Pangol1 active automatiquement le thème sombre entre 18h et 6h
    \item Réduit la fatigue visuelle lors des longues sessions
    \item Le mode peut également être activé manuellement selon vos préférences
\end{itemize}

\subsection{Fichiers Récents}
\begin{itemize}
    \item Pangol1 conserve l'historique des fichiers récemment ouverts
    \item Ces fichiers sont automatiquement rechargés au démarrage
    \item Permet une reprise rapide de votre travail
\end{itemize}

\section{Résolution de Problèmes}

\subsection{Problèmes Courants}

\subsubsection{Le fichier CSV ne s'ouvre pas}
\begin{itemize}
    \item Vérifiez que le fichier existe et n'est pas corrompu
    \item Assurez-vous que le fichier est au format CSV valide
    \item Vérifiez l'encodage du fichier (UTF-8 recommandé)
    \item Les fichiers très volumineux peuvent nécessiter plus de temps de chargement
\end{itemize}

\subsubsection{Le graphique ne s'affiche pas}
\begin{itemize}
    \item Vérifiez que vous avez sélectionné les colonnes X et Y
    \item Assurez-vous que la colonne Y contient des valeurs numériques
    \item Vérifiez que les données ne sont pas filtrées de manière à être vides
    \item Redémarrez l'application si le problème persiste
\end{itemize}

\subsubsection{Les données n'apparaissent pas correctement}
\begin{itemize}
    \item Vérifiez le délimiteur du fichier CSV (virgule ou point-virgule)
    \item Assurez-vous que la première ligne contient les en-têtes
    \item Vérifiez qu'il n'y a pas de caractères spéciaux non supportés
    \item Essayez de nettoyer les données avec un éditeur de texte
\end{itemize}

\subsection{Conseils d'Utilisation}

\begin{itemize}
    \item Préparez vos données CSV avant de les importer (supprimez les colonnes inutiles)
    \item Utilisez des noms de colonnes explicites et concis
    \item Les colonnes avec données numériques doivent contenir uniquement des nombres
    \item Évitez les valeurs manquantes ou vides dans les colonnes critiques
    \item Testez votre graphique avec un petit ensemble de données avant de traiter de gros volumes
    \item Exportez régulièrement vos graphiques importants
\end{itemize}

\section{Support}

\subsection{Obtenir de l'Aide}
\begin{itemize}
    \item \textbf{Help Menu} : Accès au manuel utilisateur
    \item \textbf{About} : Informations de version et dépendances
\end{itemize}

\subsection{Documentation Technique}
Pour les développeurs et les informations détaillées sur le projet :

\begin{itemize}
    \item \textbf{README.md} : Vue d'ensemble du projet
    \item \textbf{DOCUMENTATION.md} : Documentation technique complète
    \item \textbf{CONTRIBUTING.md} : Guide de contribution
    \item \textbf{MEMBERS.md} : Liste des contributeurs
\end{itemize}

\end{document}